\subsection{问题五分析}
前四问研究表明，禁渔的恢复收益在生态系统不同层级并非同步传递：基础食源、经济鱼类、珍稀物种及系统长期动力学行为的响应存在明显差异。据此，我们在问题五中基于统一的时间窗和指标正向化规则，选取基础食源、长江江豚、长江鲟和外来入侵压力四类核心指标，用于评估水体污染与外来物种入侵的复合生态效应。

除计算综合生态健康得分外，还需验证情景排序对赋权方法的鲁棒性。若在熵权法、CRITIC法、等权法及局部参数扰动下情景排序保持一致，则"水体污染首先削弱禁渔收益，外来物种入侵进一步恶化系统状态"的结论具有稳健性。

